\chapter{Study Two: AFI}\label{ch:study-two}

\begin{synopsis}
\section{studying indoor experience through affective interaction}

The occupant studies explore the potential identified by the scoping review through situated research in learning and working environments. They experiment with emotion as an entry point while drawing on theories and methods from Affective Interaction.

Affective Interaction provides a methodological lens for studying affect as situated, relational, and interactional.\sidenote{The \emph{affective lens} describes the broader orientation of the dissertation. \emph{Affective Interaction} is used more specifically as a methodological lens within the occupant studies.} It foregrounds how bodily responses, interpretations, technologies, spatial conditions, activities, and social circumstances contribute to experience. Applied to indoor environments, it allows HBI to ask not only whether occupants approve of a space, but how they sense it, interpret its conditions, anticipate what may occur, and make meaning from their situation.

The purpose of the occupant studies is not to determine whether a building produces positive or negative emotions. Emotions are instead used to elicit how occupants perceive, interpret, and negotiate their surroundings. An account of calm may disclose the importance of daylight, visibility, cover, social presence, or a temporary absence of demands. An account of frustration may expose a mismatch among spatial arrangements, environmental conditions, activity, and expectation.

Through these studies, indoor experience is characterised through the related micro-practices of \emph{sensing}, \emph{interpreting}, \emph{anticipating}, and \emph{making meaning}. Sensing concerns how occupants register light, sound, temperature, materiality, movement, and bodily comfort. Interpreting concerns how they understand spatial cues, environmental behaviour, technologies, and the actions of others. Anticipating concerns how expectations of interruption, visibility, movement, or environmental change shape present experience. Making meaning concerns how occupants relate these conditions to their purposes, memories, preferences, and wider circumstances.

These micro-practices are not treated as successive stages. They provide related perspectives on how indoor experience is constituted. Occupants register environmental and social conditions, assess what those conditions mean for their activity, anticipate how the situation may change, and adjust their behaviour. They decide where to sit, when to move, how long to remain, what to tolerate, and how to position themselves in relation to others.

The studies demonstrate that Affective Interaction can reveal experiential detail that may remain obscured when occupants are approached primarily as users of building systems or respondents to environmental conditions. The contribution is therefore not a method for detecting emotion. It is a methodological approach for examining how indoor environments become meaningful through relations among bodily, spatial, technological, and social conditions.
\end{synopsis}

\section{design rationale}
Explain the design space and the specific research question.

\section{procedure}
Describe participants, workshop structure, materials, prompts, and analysis.

\section{design opportunities}
\begin{enumerate}
  \item Opportunity one.
  \item Opportunity two.
  \item Opportunity three.
\end{enumerate}

\section{implications}
Translate the findings into design implications for future HBI systems.
